\chapter{Probability Theory}\label{probability}
\chapterstatus{draft}

The axioms of Kolmogorov make probability a part of measure theory: a probability is a measure of total mass $1$, a random variable is a measurable function, and its expectation is its integral. What makes
probability a subject of its own is the notion of independence, and the striking regularity that independence produces in the large: the law of large numbers, the central limit theorem, the Poisson law of rare
events, large deviations, and the universal laws of random matrices.
\begin{itemize}
\item Section~\ref{probability:section-random-variables}: random variables, distributions, distribution functions, expectation, the classical distributions.
\item Section~\ref{probability:section-independence}: independence of $\sigma$-algebras and random variables, Kolmogorov's zero-one law, sums of independent variables.
\item Section~\ref{probability:section-inequalities}: moments and the inequalities of Markov, Chebyshev, Chernoff and Hoeffding.
\item Section~\ref{probability:section-lln}: the weak and strong laws of large numbers, Kolmogorov's inequality, random series.
\item Section~\ref{probability:section-characteristic-functions}: convergence in distribution, characteristic functions, the inversion and continuity theorems.
\item Section~\ref{probability:section-clt}: the central limit theorem and the Lindeberg--Feller theorem.
\item Section~\ref{probability:section-poisson}: the Poisson approximation, infinitely divisible and stable laws.
\item Section~\ref{probability:section-large-deviations}: Cram\'er's theorem.
\item Section~\ref{probability:section-random-matrices}: Wigner's semicircle law.
\end{itemize}

\noindent\textbf{Prerequisites.} Chapter~\ref{measure-theory} (Measure and Integration), Chapter~\ref{radon-measures} (Radon Measures), Chapter~\ref{special-functions} (Special Functions) and Chapter~\ref{distributions}
(Distributions and Sobolev Spaces).

\section{Random variables and their distributions}\label{probability:section-random-variables}

\begin{definition}\label{probability:definition-probability-space}
A \emph{probability space} is a measure space $(\Omega, \mathcal{F}, P)$ with $P(\Omega) = 1$; elements of $\mathcal{F}$ are \emph{events}. A \emph{random variable} with values in a measurable space $(E, \mathcal{E})$ is a measurable
map $X \colon \Omega \to E$; a \emph{real random variable} takes values in $\RR$ with the Borel $\sigma$-algebra, and a \emph{random vector} in $\RR^d$. The \emph{distribution} or \emph{law} of $X$ is the image measure
$P_X = X_*P$ on $E$, $P_X(B) = P(X \in B)$ (Definition~\ref{measure-theory:definition-image-measure}). The \emph{distribution function} of a real random variable is $F_X(x) = P(X \leq x)$.
\end{definition}

\begin{notation}\label{probability:notation-probability}
We write $\{X \in B\}$ for $X^{-1}(B)$ and $P(X \in B)$ for its probability; ``almost surely'' (a.s.) means outside an event of probability $0$. The \emph{expectation} of a real random variable is
$E[X] = \int_\Omega X\,dP$ when it is defined, and $E[X; A] = E[X\mathbf{1}_A]$. $X \sim \mu$ means that $X$ has law $\mu$, and $X \overset{d}{=} Y$ that $X$ and $Y$ have the same law. $\operatorname{Var}(X) = E[(X - E X)^2]$ and
$\operatorname{Cov}(X, Y) = E[(X - EX)(Y - EY)]$. $N(m, \sigma^2)$ is the normal law of Example~\ref{probability:example-classical-distributions}.
\end{notation}

By the change of variables formula (Proposition~\ref{measure-theory:proposition-image-measure}), $E[g(X)] = \int_Eg\,dP_X$ for measurable $g \geq 0$ or $g$ with $g(X)$ integrable: every quantity depending only on $X$ is
computed from its law, and the underlying space $\Omega$ plays a secondary role.

\begin{proposition}\label{probability:proposition-distribution-function}
The distribution function $F$ of a real random variable is nondecreasing and right-continuous, with $F(x) \to 0$ as $x \to -\infty$ and $F(x) \to 1$ as $x \to \infty$; it determines the law, and $P(X = x) = F(x) - F(x-)$.
Conversely, every $F$ with these properties is the distribution function of a random variable: if $U$ is uniform on $(0, 1)$ and $G(u) = \inf\{x : F(x) \geq u\}$, then $G(U)$ has distribution function $F$.
\end{proposition}

\begin{proof}
Monotonicity is clear; right-continuity and the limits follow from continuity of $P$ along monotone sequences of events. $F$ determines $P_X$ on the $\pi$-system of half-lines $(-\infty, x]$, which generates the Borel sets, hence
everywhere by Dynkin's lemma (Lemma~\ref{measure-theory:lemma-dynkin}). For the converse, $G$ is nondecreasing and finite on $(0, 1)$; by right-continuity of $F$ the infimum is attained, so $G(u) \leq x$ iff $u \leq F(x)$.
Hence $P(G(U) \leq x) = P(U \leq F(x)) = F(x)$.
\end{proof}

The \emph{quantile transform} $G(U)$ of the proof is how random numbers with a given law are simulated from uniform ones.

\begin{definition}\label{probability:definition-discrete-continuous}
A real random variable is \emph{discrete} if $P(X \in S) = 1$ for a countable set $S$; its law is then determined by the \emph{probability mass function} $p(x) = P(X = x)$. It is \emph{(absolutely) continuous} with
\emph{density} $f$ if $P(X \in B) = \int_Bf\,dx$ for all Borel $B$, that is $P_X = f\,dx$; then $E[g(X)] = \int g(x)f(x)\,dx$.
\end{definition}

\begin{example}\label{probability:example-classical-distributions}
The classical laws, with mean and variance (computed from the definitions, or more easily with Proposition~\ref{probability:proposition-characteristic-function-moments}):
\begin{enumerate}
\item \emph{Bernoulli} $B(1, p)$: $P(X = 1) = p = 1 - P(X = 0)$; mean $p$, variance $p(1 - p)$. \emph{Binomial} $B(n, p)$: $P(X = k) = \binom nkp^k(1 - p)^{n-k}$, the number of successes in $n$ independent trials; mean
$np$, variance $np(1 - p)$.
\item \emph{Geometric}: $P(X = k) = (1 - p)^{k-1}p$, $k \geq 1$, the time of the first success; mean $1/p$, variance $(1 - p)/p^2$.
\item \emph{Poisson} $\mathrm{Poi}(\lambda)$: $P(X = k) = e^{-\lambda}\lambda^k/k!$; mean and variance $\lambda$.
\item \emph{Uniform} on $[a, b]$: density $\frac{1}{b - a}$ on $[a, b]$; mean $\frac{a + b}{2}$, variance $\frac{(b - a)^2}{12}$.
\item \emph{Exponential} $\mathrm{Exp}(\lambda)$: density $\lambda e^{-\lambda x}$ on $(0, \infty)$; mean $1/\lambda$, variance $1/\lambda^2$; it is \emph{memoryless}, $P(X > s + t \mid X > s) = P(X > t)$.
\item \emph{Normal} or Gaussian $N(m, \sigma^2)$: density $\frac{1}{\sigma\sqrt{2\pi}}e^{-(x - m)^2/2\sigma^2}$; mean $m$, variance $\sigma^2$. The density integrates to $1$ by Example~\ref{special-functions:example-gamma-half}.
\item \emph{Gamma} $\Gamma(\alpha, \lambda)$: density $\frac{\lambda^\alpha}{\Gamma(\alpha)}x^{\alpha-1}e^{-\lambda x}$ on $(0, \infty)$; mean $\alpha/\lambda$, variance $\alpha/\lambda^2$. For $\alpha = 1$ it is $\mathrm{Exp}(\lambda)$.
\item \emph{Cauchy}: density $\frac{1}{\pi(1 + x^2)}$. It has no mean: $E|X| = \int\frac{|x|}{\pi(1 + x^2)}dx = \infty$.
\end{enumerate}
\end{example}

\begin{counterexample}\label{probability:counterexample-law-not-variable}
The law does not determine the random variable: on $[0, 1]$ with Lebesgue measure, $X(\omega) = \omega$ and $Y(\omega) = 1 - \omega$ are both uniform, $X \overset{d}{=} Y$, but $X \neq Y$ almost everywhere. Statements about
laws (convergence in distribution, the central limit theorem) are therefore weaker than statements about the variables themselves (almost sure convergence, the strong law).
\end{counterexample}

\section{Independence}\label{probability:section-independence}

\begin{definition}\label{probability:definition-independent-sigma-algebras}
Sub-$\sigma$-algebras $\mathcal{F}_1, \mathcal{F}_2, \ldots$ of $\mathcal{F}$ are \emph{independent} if $P(A_1 \cap \cdots \cap A_m) = P(A_1)\cdots P(A_m)$ whenever $A_j \in \mathcal{F}_{k_j}$ for distinct $k_1, \ldots, k_m$.
Random variables $X_k$ are independent if the $\sigma$-algebras $\sigma(X_k) = \{X_k^{-1}(B)\}$ are, which agrees with Definition~\ref{measure-theory:definition-independent-variables}. A sequence is \emph{i.i.d.} if it is
independent and identically distributed.
\end{definition}

By Corollary~\ref{measure-theory:corollary-independent-sequences}, for any laws $\mu_k$ there is a probability space carrying independent random variables $X_k \sim \mu_k$; this is what makes the theory non-vacuous.

\begin{lemma}\label{probability:lemma-pi-system-independence}
If $\mathcal{P}_1, \ldots, \mathcal{P}_n$ are $\pi$-systems (closed under finite intersections) containing $\Omega$ with $P(A_1 \cap \cdots \cap A_n) = \prod P(A_j)$ for $A_j \in \mathcal{P}_j$, then $\sigma(\mathcal{P}_1), \ldots, \sigma(\mathcal{P}_n)$
are independent. Consequently, real random variables $X_1, \ldots, X_n$ are independent iff $P(X_1 \leq x_1, \ldots, X_n \leq x_n) = \prod_jP(X_j \leq x_j)$ for all $x_j$, iff the law of $(X_1, \ldots, X_n)$ is the product
$P_{X_1} \otimes \cdots \otimes P_{X_n}$.
\end{lemma}

\begin{proof}
Fix $A_2 \in \mathcal{P}_2, \ldots, A_n \in \mathcal{P}_n$. The sets $A_1 \in \mathcal{F}$ with $P(A_1 \cap A_2 \cap \cdots \cap A_n) = \prod P(A_j)$ form a family containing $\mathcal{P}_1$, hence $\Omega$. It is closed under countable disjoint unions by
$\sigma$-additivity, and under complements: subtract the identity for $A_1$ from the identity for $\Omega$. By Dynkin's lemma it contains $\sigma(\mathcal{P}_1)$. Repeat with each index in turn. Since each $\mathcal{P}_j$ contains $\Omega$, the
identity for all $n$ sets gives it for subfamilies. The half-lines $(-\infty, x]$ together with $\RR$ form a $\pi$-system generating the Borel sets, which gives the second statement; the third follows because the product
measure is characterised by its values on boxes (Lemma~\ref{measure-theory:lemma-dynkin}).
\end{proof}

\begin{corollary}\label{probability:corollary-independence-products}
\begin{enumerate}
\item (Grouping) If $X_1, X_2, \ldots$ are independent, then $f(X_1, \ldots, X_n)$ and $g(X_{n+1}, X_{n+2}, \ldots)$ are independent for measurable $f$, $g$; more generally, disjoint groups of independent variables generate
independent $\sigma$-algebras.
\item If $X$, $Y$ are independent and integrable, then $XY$ is integrable and $E[XY] = E[X]E[Y]$; so $\operatorname{Cov}(X, Y) = 0$ and $\operatorname{Var}(X + Y) = \operatorname{Var}X + \operatorname{Var}Y$ when $X, Y \in L^2$.
\item If $X$, $Y$ are independent real random variables, $P_{X+Y} = P_X * P_Y$, where $(\mu * \nu)(B) = \int\nu(B - x)\,d\mu(x)$; if they have densities $f$, $g$, then $X + Y$ has density $f * g$.
\end{enumerate}
\end{corollary}

\begin{proof}
(1) The events $\{X_1 \in B_1, \ldots, X_n \in B_n\}$, and those $\{X_{n+1} \in B_{n+1}, \ldots, X_m \in B_m\}$ for all $m$, form $\pi$-systems containing $\Omega$ that generate $\sigma(X_1, \ldots, X_n)$ and $\sigma(X_{n+1}, \ldots)$, and
they satisfy the product rule; apply Lemma~\ref{probability:lemma-pi-system-independence}. (2) By the lemma and Fubini's theorem (Theorem~\ref{measure-theory:theorem-fubini}),
$E|XY| = \int|x||y|\,d(P_X \otimes P_Y) = E|X|E|Y|$, and the same computation without absolute values. (3) $P(X + Y \in B) = \int\int\mathbf{1}_B(x + y)\,dP_Y(y)\,dP_X(x)$ by Fubini; with densities, substitute $y = z - x$.
\end{proof}

\begin{example}\label{probability:example-sums}
If $X \sim \mathrm{Poi}(\lambda)$ and $Y \sim \mathrm{Poi}(\mu)$ are independent, $P(X + Y = n) = \sum_k e^{-\lambda}\frac{\lambda^k}{k!}e^{-\mu}\frac{\mu^{n-k}}{(n - k)!} = e^{-\lambda-\mu}\frac{(\lambda + \mu)^n}{n!}$ by the binomial theorem:
$X + Y \sim \mathrm{Poi}(\lambda + \mu)$. A sum of $n$ independent $\mathrm{Exp}(\lambda)$ variables is $\Gamma(n, \lambda)$, by induction on the convolution $\int_0^x\frac{\lambda^{n-1}y^{n-2}e^{-\lambda y}}{(n - 2)!}\lambda e^{-\lambda(x - y)}dy$.
Sums of independent normals are normal (Example~\ref{probability:example-characteristic-functions}).
\end{example}

\begin{counterexample}\label{probability:counterexample-uncorrelated}
Uncorrelated does not imply independent: if $X \sim N(0, 1)$ and $Y = X^2$, then $\operatorname{Cov}(X, Y) = E[X^3] = 0$, but $P(|X| \leq 1, Y > 1) = 0 \neq P(|X| \leq 1)P(Y > 1)$. And pairwise independence does not imply
independence (Example~\ref{measure-theory:example-pairwise-independent}).
\end{counterexample}

\begin{definition}\label{probability:definition-tail}
For random variables $X_1, X_2, \ldots$, the \emph{tail $\sigma$-algebra} is $\mathcal{T} = \bigcap_n\sigma(X_n, X_{n+1}, \ldots)$. Its events do not depend on any finite number of the $X_k$: for real variables, $\{\sum X_k$
converges$\}$ and $\{\limsup_n\frac1n(X_1 + \cdots + X_n) \leq c\}$ are tail events.
\end{definition}

\begin{theorem}[Kolmogorov's zero-one law]\label{probability:theorem-zero-one}
If $X_1, X_2, \ldots$ are independent, every tail event has probability $0$ or $1$, and every $\mathcal{T}$-measurable random variable with values in $[-\infty, \infty]$ is almost surely constant.
\end{theorem}

\begin{proof}
By Corollary~\ref{probability:corollary-independence-products}(1), $\sigma(X_1, \ldots, X_n)$ is independent of $\sigma(X_{n+1}, \ldots) \supseteq \mathcal{T}$. So $\mathcal{T}$ is independent of the $\pi$-system
$\bigcup_n\sigma(X_1, \ldots, X_n)$, hence by Lemma~\ref{probability:lemma-pi-system-independence} of the $\sigma$-algebra it generates, which contains $\mathcal{T}$. So every $A \in \mathcal{T}$ is independent of itself:
$P(A) = P(A)^2$. For a $\mathcal{T}$-measurable $Y$, $P(Y \leq c) \in \{0, 1\}$ for all $c$, so $Y = \sup\{c : P(Y \leq c) = 0\}$ almost surely.
\end{proof}

So a random series of independent terms converges with probability $0$ or $1$, and $\limsup S_n/n$ is almost surely constant. The second Borel--Cantelli lemma (Theorem~\ref{measure-theory:theorem-borel-cantelli}(2)) is a
quantitative instance.

\section{Moments and inequalities}\label{probability:section-inequalities}

\begin{definition}\label{probability:definition-moments}
The \emph{$k$-th moment} of $X$ is $E[X^k]$ when $E|X|^k < \infty$. The \emph{moment generating function} is $M_X(\lambda) = E[e^{\lambda X}] \in (0, \infty]$, and $\Lambda_X = \log M_X$ is the \emph{cumulant generating function}.
\end{definition}

The basic inequalities are those of measure theory read probabilistically: Markov's $P(|X| \geq t) \leq E|X|/t$, Chebyshev's $P(|X - EX| \geq t) \leq \operatorname{Var}X/t^2$ (Remark~\ref{measure-theory:remark-markov-chebyshev}),
Jensen's $\varphi(EX) \leq E\varphi(X)$ for convex $\varphi$ (Proposition~\ref{measure-theory:proposition-jensen}), and H\"older's inequality, which gives $\|X\|_p \leq \|X\|_q$ for $p \leq q$ on a probability space. Applying
Markov's inequality to $e^{\lambda X}$ gives exponential bounds.

\begin{proposition}[Chernoff bound]\label{probability:proposition-chernoff}
For every real $X$ and $t$, $P(X \geq t) \leq \inf_{\lambda \geq 0}e^{-\lambda t}M_X(\lambda)$. If $S_n$ is a sum of $n$ i.i.d.\ copies of $X$, $P(S_n \geq nt) \leq \exp\bigl(-n\sup_{\lambda \geq 0}(\lambda t - \Lambda_X(\lambda))\bigr)$.
\end{proposition}

\begin{proof}
$P(X \geq t) \leq P(e^{\lambda X} \geq e^{\lambda t}) \leq e^{-\lambda t}E[e^{\lambda X}]$ by Markov, for $\lambda \geq 0$. For $S_n$, $M_{S_n} = M_X^n$ by Corollary~\ref{probability:corollary-independence-products}(2).
\end{proof}

\begin{lemma}[Hoeffding]\label{probability:lemma-hoeffding}
If $a \leq X \leq b$ and $E[X] = 0$, then $E[e^{\lambda X}] \leq e^{\lambda^2(b - a)^2/8}$ for all real $\lambda$.
\end{lemma}

\begin{proof}
We may assume $a < 0 < b$ (if $a = 0$ or $b = 0$ then $X = 0$). By convexity, $e^{\lambda x} \leq \frac{b - x}{b - a}e^{\lambda a} + \frac{x - a}{b - a}e^{\lambda b}$ on $[a, b]$; taking expectations,
$E[e^{\lambda X}] \leq \frac{b}{b - a}e^{\lambda a} - \frac{a}{b - a}e^{\lambda b} = e^{L(h)}$ with $h = \lambda(b - a)$, $p = \frac{-a}{b - a} \in (0, 1)$ and $L(h) = -hp + \log(1 - p + pe^h)$. Then $L(0) = L^{\prime}(0) = 0$ and
$L^{\prime\prime}(h) = q(1 - q) \leq \frac14$ with $q = \frac{pe^h}{1 - p + pe^h}$, so $L(h) \leq \frac{h^2}{8}$ by Taylor's theorem.
\end{proof}

\begin{theorem}[Hoeffding's inequality]\label{probability:theorem-hoeffding}
Let $X_1, \ldots, X_n$ be independent with $a_i \leq X_i \leq b_i$, and $S = \sum X_i$. Then for $t > 0$,
\[
P(S - ES \geq t) \leq \exp\Bigl(-\frac{2t^2}{\sum_i(b_i - a_i)^2}\Bigr).
\]
\end{theorem}

\begin{proof}
By Lemma~\ref{probability:lemma-hoeffding} applied to $X_i - EX_i$ and independence, $E[e^{\lambda(S - ES)}] \leq e^{\lambda^2D/8}$ with $D = \sum(b_i - a_i)^2$. The Chernoff bound gives $e^{-\lambda t + \lambda^2D/8}$, minimal at $\lambda = 4t/D$.
\end{proof}

\begin{example}\label{probability:example-hoeffding}
For $n$ fair coin tosses with $S$ heads, $P(|S - \frac n2| \geq t) \leq 2e^{-2t^2/n}$; with $t = 0.1n$, the probability that the proportion of heads deviates from $\frac12$ by $0.1$ is at most $2e^{-0.02n}$, about $4 \cdot 10^{-9}$
for $n = 1000$. Chebyshev's inequality gives only $\frac{n/4}{(0.1n)^2} = 0.025$. Concentration inequalities of this kind are the basic tool of the probabilistic method (Chapter~\ref{probabilistic-method}) and of statistical
learning theory.
\end{example}

\section{Laws of large numbers}\label{probability:section-lln}

\begin{theorem}[Weak law of large numbers]\label{probability:theorem-weak-law}
Let $X_1, X_2, \ldots$ be pairwise uncorrelated with $E[X_k] = m$ and $\operatorname{Var}X_k \leq C$. Then $S_n/n \to m$ in probability: $P(|S_n/n - m| \geq \varepsilon) \leq \frac{C}{n\varepsilon^2}$.
\end{theorem}

\begin{proof}
$\operatorname{Var}(S_n) = \sum\operatorname{Var}X_k \leq nC$, and Chebyshev's inequality gives the bound.
\end{proof}

\begin{example}[Bernstein polynomials]\label{probability:example-bernstein}
Let $f \in C[0, 1]$ and $B_nf(x) = \sum_{k=0}^nf(\frac kn)\binom nkx^k(1 - x)^{n-k} = E[f(S_n/n)]$, with $S_n \sim B(n, x)$. Given $\varepsilon$, choose $\delta$ with $|f(y) - f(x)| < \varepsilon$ for $|y - x| < \delta$. Then
$|B_nf(x) - f(x)| \leq E|f(S_n/n) - f(x)| \leq \varepsilon + 2\|f\|_\infty P(|S_n/n - x| \geq \delta) \leq \varepsilon + \frac{2\|f\|_\infty x(1 - x)}{n\delta^2} \leq \varepsilon + \frac{\|f\|_\infty}{2n\delta^2}$, uniformly in $x$. So the polynomials $B_nf$
converge uniformly to $f$: a probabilistic proof of Weierstrass' approximation theorem (Corollary~\ref{real-analysis:corollary-weierstrass}).
\end{example}

\begin{theorem}[Strong law of large numbers; Etemadi]\label{probability:theorem-strong-law}
Let $X_1, X_2, \ldots$ be pairwise independent, identically distributed and integrable, with $E[X_1] = m$. Then $S_n/n \to m$ almost surely.
\end{theorem}

\begin{proof}
\emph{Reduction to $X_k \geq 0$.} The positive and negative parts $X_k^\pm$ are again pairwise independent and identically distributed, so we may assume $X_k \geq 0$.
\emph{Truncation.} Let $Y_k = X_k\mathbf{1}_{\{X_k \leq k\}}$ and $T_n = Y_1 + \cdots + Y_n$. Since $\sum_kP(X_k > k) \leq \int_0^\infty P(X_1 > t)\,dt = E[X_1] < \infty$ (layer cake formula,
Proposition~\ref{measure-theory:proposition-layer-cake}), Borel--Cantelli shows that almost surely $X_k = Y_k$ for all large $k$; so it suffices to prove $T_n/n \to m$.
\emph{A variance bound.} By the layer cake formula $E[Y_k^2] = \int_0^\infty2yP(Y_k > y)\,dy \leq \int_0^k2yP(X_1 > y)\,dy$. For $y \geq 0$, $2y\sum_{k > y}k^{-2} \leq 4$ (for $y < 1$ use $\sum k^{-2} \leq 2$; for $y \geq 1$,
$\sum_{k > y}k^{-2} \leq \frac{1}{\lfloor y\rfloor} \leq \frac2y$). Hence, by Tonelli,
\[
\sum_k\frac{E[Y_k^2]}{k^2} \leq \int_0^\infty\Bigl(2y\sum_{k > y}k^{-2}\Bigr)P(X_1 > y)\,dy \leq 4E[X_1].
\]
\emph{Along a geometric subsequence.} Fix $\alpha > 1$ and let $k(n) = \lfloor\alpha^n\rfloor$. The $Y_k$ are pairwise independent, so $\operatorname{Var}T_{k(n)} = \sum_{j \leq k(n)}\operatorname{Var}Y_j \leq \sum_{j \leq k(n)}E[Y_j^2]$. By Chebyshev,
\[
\begin{aligned}
\sum_nP\bigl(|T_{k(n)} - E T_{k(n)}| > \varepsilon k(n)\bigr) &\leq \frac{1}{\varepsilon^2}\sum_n\frac{1}{k(n)^2}\sum_{j \leq k(n)}E[Y_j^2] \\
&= \frac{1}{\varepsilon^2}\sum_jE[Y_j^2]\sum_{n : k(n) \geq j}\frac{1}{k(n)^2}.
\end{aligned}
\]
As $k(n) \geq \alpha^n/2$, the inner sum is at most $4\sum_{n : \alpha^n \geq j}\alpha^{-2n} \leq \frac{4}{1 - \alpha^{-2}}j^{-2}$, so the whole sum is finite by the variance bound. By Borel--Cantelli and letting $\varepsilon \to 0$ along a sequence,
$(T_{k(n)} - ET_{k(n)})/k(n) \to 0$ almost surely. Since $E[Y_j] \to E[X_1] = m$ by monotone convergence, $ET_{k(n)}/k(n) \to m$ (Ces\`aro), and so $T_{k(n)}/k(n) \to m$ almost surely.
\emph{Interpolation.} For $k(n) \leq j < k(n + 1)$, as $T$ is nondecreasing, $\frac{T_{k(n)}}{k(n + 1)} \leq \frac{T_j}{j} \leq \frac{T_{k(n+1)}}{k(n)}$. Since $k(n + 1)/k(n) \to \alpha$, $\frac m\alpha \leq \liminf\frac{T_j}{j} \leq \limsup\frac{T_j}{j} \leq \alpha m$
almost surely. Let $\alpha \downarrow 1$ along a sequence.
\end{proof}

This proof, which needs only pairwise independence and no fourth moments, is due to Etemadi \cite{etemadi-1981}; for the classical approach via Kolmogorov's inequality see \cite{billingsley-probability}.

\begin{example}\label{probability:example-strong-law}
\begin{enumerate}
\item \emph{Frequencies.} For i.i.d.\ $X_k$ and a Borel set $B$, $\frac1n\#\{k \leq n : X_k \in B\} \to P(X_1 \in B)$ almost surely: probability is limiting frequency, which is the frequentist's definition turned into a theorem. With
the binary digits of a uniform number this gives Borel's theorem on normal numbers (Theorem~\ref{measure-theory:theorem-borel-normal}).
\item \emph{Monte Carlo integration.} For $U_k$ i.i.d.\ uniform on $[0, 1]^d$ and $f \in L^1$, $\frac1n\sum_{k \leq n}f(U_k) \to \int f$ almost surely, with error of order $n^{-1/2}$ (by the central limit theorem) in every dimension.
\end{enumerate}
\end{example}

\begin{counterexample}\label{probability:counterexample-no-mean}
Integrability is necessary: if $X_k$ are i.i.d.\ with $E|X_1| = \infty$, then $\limsup_n|S_n|/n = \infty$ almost surely. Indeed, for every $c > 0$, $\sum_nP(|X_n| > cn) \geq \int_1^\infty P(|X_1| > ct)\,dt = \infty$, so by the second
Borel--Cantelli lemma $|X_n| > cn$ infinitely often; since $|X_n| \leq |S_n| + |S_{n-1}|$, $\max(|S_n|/n, |S_{n-1}|/(n - 1)) > c/2$ infinitely often. For the Cauchy law $S_n/n$ is again Cauchy for every $n$
(Example~\ref{probability:example-characteristic-functions}), so averaging does not help at all.
\end{counterexample}

\begin{theorem}[Kolmogorov's maximal inequality]\label{probability:theorem-kolmogorov-inequality}
Let $X_1, \ldots, X_n$ be independent with mean $0$ and finite variance, and $S_k = X_1 + \cdots + X_k$. Then for $\lambda > 0$,
$P(\max_{k \leq n}|S_k| \geq \lambda) \leq \operatorname{Var}(S_n)/\lambda^2$.
\end{theorem}

\begin{proof}
Let $A_k = \{|S_k| \geq \lambda, |S_j| < \lambda \text{ for } j < k\}$, disjoint with union $\{\max|S_k| \geq \lambda\}$. Then $E[S_n^2] \geq \sum_kE[S_n^2; A_k] \geq \sum_kE[S_k^2 + 2S_k(S_n - S_k); A_k]$. The variable $S_k\mathbf{1}_{A_k}$ depends
on $X_1, \ldots, X_k$, hence is independent of $S_n - S_k$, which has mean $0$; so the cross terms vanish and $E[S_n^2] \geq \sum_kE[S_k^2; A_k] \geq \lambda^2\sum_kP(A_k)$.
\end{proof}

\begin{theorem}[Random series]\label{probability:theorem-random-series}
\begin{enumerate}
\item If $X_k$ are independent with mean $0$ and $\sum_k\operatorname{Var}X_k < \infty$, then $\sum_kX_k$ converges almost surely.
\item (Kolmogorov's three-series theorem, sufficiency) Let $X_k$ be independent, $A > 0$ and $Y_k = X_k\mathbf{1}_{\{|X_k| \leq A\}}$. If $\sum_kP(|X_k| > A)$, $\sum_kE[Y_k]$ and $\sum_k\operatorname{Var}Y_k$ converge, then $\sum X_k$
converges almost surely.
\end{enumerate}
\end{theorem}

\begin{proof}
(1) By Theorem~\ref{probability:theorem-kolmogorov-inequality} applied to $X_{m+1}, \ldots, X_n$ and letting $n \to \infty$, $P(\sup_{k > m}|S_k - S_m| \geq \varepsilon) \leq \varepsilon^{-2}\sum_{k > m}\operatorname{Var}X_k \to 0$. So
$\sup_{j,k > m}|S_j - S_k|$ tends to $0$ in probability; being decreasing in $m$, it tends to $0$ almost surely, and $(S_n)$ is almost surely Cauchy. (2) By (1), $\sum(Y_k - EY_k)$ converges almost surely, hence so does $\sum Y_k$;
by Borel--Cantelli $X_k = Y_k$ for all large $k$ almost surely.
\end{proof}

The three-series theorem is also necessary: if $\sum X_k$ converges with positive probability, the three series converge for every $A > 0$ \cite{durrett-probability}.

\begin{example}\label{probability:example-random-harmonic}
Let $\varepsilon_k = \pm1$ be independent fair signs. Then $\sum_k\varepsilon_k/k^\alpha$ converges almost surely for $\alpha > \frac12$, by Theorem~\ref{probability:theorem-random-series}(1), since $\sum k^{-2\alpha} < \infty$; for
$\alpha \leq \frac12$ it diverges almost surely (by the necessity of the three series). The random harmonic series $\sum\pm\frac1k$ converges with probability $1$, although $\sum\frac1k$ diverges.
\end{example}

\section{Convergence in distribution and characteristic functions}\label{probability:section-characteristic-functions}

\begin{definition}\label{probability:definition-convergence-distribution}
Real random variables $X_n$ \emph{converge in distribution} to $X$, written $X_n \Rightarrow X$, if $E[h(X_n)] \to E[h(X)]$ for every bounded continuous $h \colon \RR \to \RR$; that is, the laws $P_{X_n}$ converge narrowly to $P_X$
(Definition~\ref{radon-measures:definition-narrow}). The variables may live on different probability spaces. The \emph{characteristic function} of $X$ is $\varphi_X(t) = E[e^{itX}] = \int e^{itx}\,dP_X(x)$, $t \in \RR$.
\end{definition}

\begin{notation}\label{probability:notation-characteristic}
$\Rightarrow$ denotes convergence in distribution, also written $X_n \Rightarrow \mu$ for a law $\mu$. The characteristic function is a Fourier transform with the probabilists' normalisation: in the notation of
Chapter~\ref{distributions}, $\varphi_X(t) = \widehat{P_X}(-t/2\pi)$.
\end{notation}

\begin{proposition}\label{probability:proposition-convergence-distribution}
\begin{enumerate}
\item Almost sure convergence implies convergence in probability, which implies convergence in distribution.
\item $X_n \Rightarrow X$ iff $F_{X_n}(x) \to F_X(x)$ at every point $x$ where $F_X$ is continuous.
\item If $X_n \Rightarrow c$ for a constant $c$, then $X_n \to c$ in probability.
\end{enumerate}
\end{proposition}

\begin{proof}
(1) The first implication is Proposition~\ref{measure-theory:proposition-modes}. If $X_n \to X$ in probability and $h$ is bounded continuous, every subsequence has a further subsequence converging almost surely
(Theorem~\ref{measure-theory:theorem-cauchy-in-measure}), along which $E[h(X_n)] \to E[h(X)]$ by dominated convergence; so the whole sequence converges.
(2) If $X_n \Rightarrow X$ and $\delta > 0$, let $h$ be continuous, equal to $1$ on $(-\infty, x]$, $0$ on $[x + \delta, \infty)$, linear in between. Then $\limsup F_{X_n}(x) \leq \lim E[h(X_n)] = E[h(X)] \leq F_X(x + \delta)$, and similarly
$\liminf F_{X_n}(x) \geq F_X(x - \delta)$; at a continuity point this gives convergence. Conversely, suppose $F_n = F_{X_n} \to F = F_X$ at continuity points, and let $G_n$, $G$ be the quantile functions of
Proposition~\ref{probability:proposition-distribution-function}. If $G$ is continuous at $u \in (0, 1)$ and $y < G(u) < z$ are continuity points of $F$, then $F(y) < u$ (as $G(u) \leq y$ iff $u \leq F(y)$), so $F_n(y) < u$ and $G_n(u) > y$ for large $n$.
Also $F(z) > u$: by continuity of $G$ at $u$ there is $u^{\prime} \in (u, 1)$ with $G(u^{\prime}) < z$, and then $u^{\prime} \leq F(G(u^{\prime})) \leq F(z)$. So $F_n(z) > u$, hence $G_n(u) \leq z$, for large $n$. Hence $G_n(u) \to G(u)$ at the continuity points of $G$, which are all but countably many. With $U$ uniform, $G_n(U) \overset{d}{=} X_n$ converges almost surely to $G(U) \overset{d}{=} X$, and (1) applies.
(3) $P(|X_n - c| > \varepsilon) \leq 1 - F_n(c + \varepsilon) + F_n(c - \varepsilon) \to 0$ by (2).
\end{proof}

\begin{counterexample}\label{probability:counterexample-distribution-not-probability}
Convergence in distribution does not imply convergence in probability: if $X \sim N(0, 1)$ and $X_n = -X$ for all $n$, then $X_n \Rightarrow X$ (same law), but $|X_n - X| = 2|X|$ does not tend to $0$. The convergence $F_n(x) \to F(x)$
may fail at discontinuities of $F$: $X_n = \frac1n \Rightarrow 0$, but $F_{X_n}(0) = 0 \not\to 1 = F_0(0)$.
\end{counterexample}

\begin{proposition}\label{probability:proposition-characteristic-function-moments}
$\varphi_X$ is uniformly continuous, $\varphi_X(0) = 1$, $|\varphi_X| \leq 1$, $\varphi_{aX+b}(t) = e^{itb}\varphi_X(at)$, $\varphi_{-X} = \overline{\varphi_X}$, and $\varphi_{X+Y} = \varphi_X\varphi_Y$ for independent $X$, $Y$. If $E|X|^k < \infty$, then
$\varphi_X$ is $k$ times continuously differentiable with $\varphi_X^{(j)}(0) = i^jE[X^j]$ for $j \leq k$; in particular, if $E[X] = 0$ and $E[X^2] = \sigma^2$, then $\varphi_X(t) = 1 - \frac{\sigma^2t^2}{2} + o(t^2)$ as $t \to 0$.
\end{proposition}

\begin{proof}
$|\varphi(t + h) - \varphi(t)| \leq E|e^{ihX} - 1|$, which tends to $0$ with $h$ by dominated convergence, uniformly in $t$. Multiplicativity is Corollary~\ref{probability:corollary-independence-products}(2) for the complex
variables $e^{itX}$, $e^{itY}$ (apply it to real and imaginary parts). Differentiation under the integral sign is justified by $|\frac{d^j}{dt^j}e^{itx}| = |x|^j$ and dominated convergence. The last statement is Taylor's formula
at order $2$.
\end{proof}

\begin{example}\label{probability:example-characteristic-functions}
\begin{enumerate}
\item $\mathrm{Poi}(\lambda)$: $\varphi(t) = \sum_ke^{-\lambda}\frac{(\lambda e^{it})^k}{k!} = e^{\lambda(e^{it} - 1)}$. $B(n, p)$: $(1 - p + pe^{it})^n$. Uniform on $[-1, 1]$: $\frac{\sin t}{t}$.
\item $N(0, 1)$: $\varphi(t) = e^{-t^2/2}$, by Theorem~\ref{distributions:theorem-fourier-schwartz}(1) after the substitution $x = \sqrt{2\pi}y$; so $N(m, \sigma^2)$ has $\varphi(t) = e^{imt - \sigma^2t^2/2}$, and by multiplicativity and
uniqueness (Theorem~\ref{probability:theorem-inversion}) independent $N(m_1, \sigma_1^2)$ and $N(m_2, \sigma_2^2)$ variables have an $N(m_1 + m_2, \sigma_1^2 + \sigma_2^2)$ sum.
\item Cauchy: $\varphi(t) = e^{-|t|}$. Indeed $\int e^{-|t|}e^{-itx}dt = \frac{2}{1 + x^2}$ by direct integration, so by the inversion formula of Theorem~\ref{probability:theorem-inversion}, the law with characteristic function $e^{-|t|}$ has
density $\frac{1}{2\pi}\cdot\frac{2}{1 + x^2}$. Consequently, for i.i.d.\ Cauchy $X_k$, $\varphi_{S_n/n}(t) = (e^{-|t|/n})^n = e^{-|t|}$: $S_n/n$ is Cauchy.
\end{enumerate}
\end{example}

\begin{theorem}[Inversion and uniqueness]\label{probability:theorem-inversion}
A probability measure $\mu$ on $\RR$ is determined by its characteristic function $\varphi$. If $\varphi \in L^1(\RR)$, then $\mu$ has the bounded continuous density $f(x) = \frac{1}{2\pi}\int e^{-itx}\varphi(t)\,dt$.
\end{theorem}

\begin{proof}
For $\sigma > 0$ let $g_\sigma(u) = \frac{1}{\sigma\sqrt{2\pi}}e^{-u^2/2\sigma^2}$, the density of $N(0, \sigma^2)$. By Example~\ref{probability:example-characteristic-functions}(2) with the roles of variables exchanged,
$g_\sigma(u) = \frac{1}{2\pi}\int e^{itu}e^{-\sigma^2t^2/2}dt$. The measure $\mu * g_\sigma$ (the law of $X + \sigma Z$ with $X \sim \mu$ and $Z \sim N(0, 1)$ independent) has density
\[
\begin{aligned}
f_\sigma(x) = \int g_\sigma(x - y)\,d\mu(y) &= \frac{1}{2\pi}\int\int e^{it(x - y)}e^{-\sigma^2t^2/2}\,dt\,d\mu(y) \\
&= \frac{1}{2\pi}\int e^{-itx}\varphi(t)e^{-\sigma^2t^2/2}\,dt,
\end{aligned}
\]
by Fubini (the integrand is integrable) and the substitution $t \mapsto -t$. So $f_\sigma$, hence $\mu * g_\sigma$, is determined by $\varphi$. As $\sigma \to 0$, $X + \sigma Z \to X$ almost surely, so $\mu * g_\sigma \to \mu$ narrowly
(Proposition~\ref{probability:proposition-convergence-distribution}), and $\mu$ is determined, since narrow limits are unique (two finite Radon measures with the same integrals on $C_c$ coincide, by Riesz--Markov). If
$\varphi \in L^1$, $f_\sigma \to f$ uniformly by dominated convergence, so $\int hf\,dx = \lim\int hf_\sigma\,dx = \int h\,d\mu$ for $h \in C_c(\RR)$, and $\mu = f\,dx$.
\end{proof}

\begin{theorem}[L\'evy's continuity theorem]\label{probability:theorem-levy-continuity}
Let $\mu_n$ be probability measures on $\RR$ with characteristic functions $\varphi_n$. If $\mu_n \to \mu$ narrowly, then $\varphi_n \to \varphi_\mu$ pointwise. Conversely, if $\varphi_n \to \varphi$ pointwise and $\varphi$ is continuous at $0$,
then $\varphi$ is the characteristic function of a probability measure $\mu$ and $\mu_n \to \mu$ narrowly.
\end{theorem}

\begin{proof}
The first statement holds because $x \mapsto e^{itx}$ is bounded and continuous. \emph{A tightness estimate.} For a probability measure $\nu$ with characteristic function $\psi$ and $u > 0$, Fubini gives
$\frac1u\int_{-u}^u(1 - \psi(t))\,dt = \int\bigl(2 - \frac{2\sin ux}{ux}\bigr)d\nu(x) \geq \nu(\{|x| \geq 2/u\})$, since the integrand is $\geq 0$, and $\geq 1$ when $|ux| \geq 2$ because then $|\frac{\sin ux}{ux}| \leq \frac12$.
\emph{Tightness.} Given $\varepsilon$, choose $u$ with $\frac1u\int_{-u}^u|1 - \varphi| < \varepsilon$, by continuity at $0$ and $\varphi(0) = \lim\varphi_n(0) = 1$. By dominated convergence $\frac1u\int_{-u}^u(1 - \varphi_n) \to \frac1u\int_{-u}^u(1 - \varphi)$, so
$\mu_n(\{|x| \geq 2/u\}) < 2\varepsilon$ for large $n$; enlarging the interval for the finitely many remaining $n$, $(\mu_n)$ is tight.
\emph{Conclusion.} By Prokhorov's theorem (Theorem~\ref{radon-measures:theorem-prokhorov}) every subsequence has a further subsequence converging narrowly to some $\nu$, whose characteristic function is
$\lim\varphi_n = \varphi$ by the first part. By Theorem~\ref{probability:theorem-inversion}, all these limits are the same measure $\mu$, and then the whole sequence converges to $\mu$: otherwise some bounded continuous $h$ and
some subsequence would keep $|\int h\,d\mu_n - \int h\,d\mu| \geq \varepsilon$, but that subsequence has a further subsequence converging to $\mu$.
\end{proof}

\begin{counterexample}\label{probability:counterexample-continuity-at-zero}
Continuity of the limit at $0$ is needed: for $\mu_n = N(0, n)$, $\varphi_n(t) = e^{-nt^2/2} \to \mathbf{1}_{\{0\}}(t)$, which is not continuous at $0$, and indeed $\mu_n$ does not converge narrowly: the mass escapes to infinity.
\end{counterexample}

\section{The central limit theorem}\label{probability:section-clt}

\begin{lemma}\label{probability:lemma-products}
\begin{enumerate}
\item If $z_1, \ldots, z_n, w_1, \ldots, w_n$ are complex numbers of modulus $\leq 1$, then $|\prod z_k - \prod w_k| \leq \sum|z_k - w_k|$.
\item For real $x$, $|e^{ix} - 1 - ix + \frac{x^2}{2}| \leq \min(\frac{|x|^3}{6}, x^2)$.
\end{enumerate}
\end{lemma}

\begin{proof}
(1) Induction, using $z_1Z - w_1W = z_1(Z - W) + (z_1 - w_1)W$. (2) Taylor's formula with integral remainder: the left side is $|\frac{i^3}{2}\int_0^x(x - s)^2e^{is}ds| \leq \frac{|x|^3}{6}$. Also
$|e^{ix} - 1 - ix| = |\int_0^xi(e^{is} - 1)\,ds| \leq |\int_0^x|s|\,ds| = \frac{x^2}{2}$, using $|e^{is} - 1| \leq |s|$; adding $\frac{x^2}{2}$ gives the bound $x^2$.
\end{proof}

\begin{theorem}[Lindeberg--Feller]\label{probability:theorem-lindeberg}
For each $n$ let $X_{n,1}, \ldots, X_{n,n}$ be independent with mean $0$, and suppose
\begin{enumerate}
\item $\sum_mE[X_{n,m}^2] \to \sigma^2 > 0$;
\item (\emph{Lindeberg's condition}) $\sum_mE[X_{n,m}^2; |X_{n,m}| > \varepsilon] \to 0$ for every $\varepsilon > 0$.
\end{enumerate}
Then $S_n = X_{n,1} + \cdots + X_{n,n} \Rightarrow N(0, \sigma^2)$.
\end{theorem}

\begin{proof}
Let $\sigma_{n,m}^2 = E[X_{n,m}^2]$ and $\varphi_{n,m}$ the characteristic functions. Since $\sigma_{n,m}^2 \leq \varepsilon^2 + E[X_{n,m}^2; |X_{n,m}| > \varepsilon]$, (2) gives $\max_m\sigma_{n,m}^2 \to 0$. Fix $t$. By
Lemma~\ref{probability:lemma-products}(2), splitting according to $|X_{n,m}| \leq \varepsilon$,
\[
\begin{aligned}
\bigl|\varphi_{n,m}(t) - (1 - \tfrac12t^2\sigma_{n,m}^2)\bigr| &\leq E\min(|tX_{n,m}|^3, |tX_{n,m}|^2) \\
&\leq \varepsilon|t|^3\sigma_{n,m}^2 + t^2E[X_{n,m}^2; |X_{n,m}| > \varepsilon].
\end{aligned}
\]
For large $n$ all $c_{n,m} = \frac12t^2\sigma_{n,m}^2$ lie in $[0, \frac12]$, so Lemma~\ref{probability:lemma-products}(1) and (1), (2) give
$\limsup_n|\prod_m\varphi_{n,m}(t) - \prod_m(1 - c_{n,m})| \leq \varepsilon|t|^3\sigma^2$, for every $\varepsilon$. Finally, $|\log(1 - c) + c| \leq c^2$ for $0 \leq c \leq \frac12$, so
$|\sum_m\log(1 - c_{n,m}) + \sum_mc_{n,m}| \leq \max_mc_{n,m}\sum_mc_{n,m} \to 0$, and $\prod_m(1 - c_{n,m}) \to e^{-t^2\sigma^2/2}$. So $\varphi_{S_n}(t) \to e^{-\sigma^2t^2/2}$, and L\'evy's theorem concludes.
\end{proof}

\begin{corollary}[Central limit theorem]\label{probability:corollary-clt}
If $X_1, X_2, \ldots$ are i.i.d.\ with mean $m$ and variance $\sigma^2 \in (0, \infty)$, then $\frac{S_n - nm}{\sigma\sqrt n} \Rightarrow N(0, 1)$, that is,
$P\bigl(\frac{S_n - nm}{\sigma\sqrt n} \leq x\bigr) \to \frac{1}{\sqrt{2\pi}}\int_{-\infty}^xe^{-u^2/2}du$ for every $x$.
\end{corollary}

\begin{proof}
Apply Theorem~\ref{probability:theorem-lindeberg} to $X_{n,m} = (X_m - m)/(\sigma\sqrt n)$: $\sum_mE[X_{n,m}^2] = 1$, and the Lindeberg sum is $\sigma^{-2}E[(X_1 - m)^2; |X_1 - m| > \varepsilon\sigma\sqrt n] \to 0$ by dominated
convergence. The distribution function of $N(0, 1)$ is continuous, so Proposition~\ref{probability:proposition-convergence-distribution}(2) gives the second form.
\end{proof}

\begin{example}[De Moivre--Laplace]\label{probability:example-de-moivre}
For $S_n \sim B(n, p)$, $\frac{S_n - np}{\sqrt{np(1 - p)}} \Rightarrow N(0, 1)$. With $n = 10000$ fair coin tosses, $P(4900 \leq S_n \leq 5100) \approx P(|Z| \leq 2) \approx 0.954$. De Moivre found this approximation in 1733 for
$p = \frac12$, using Stirling's formula on the binomial coefficients.
\end{example}

\begin{remark}\label{probability:remark-clt}
\begin{enumerate}
\item (\emph{Berry--Esseen}) If moreover $E|X_1|^3 = \rho < \infty$, then $\sup_x|P(\frac{S_n - nm}{\sigma\sqrt n} \leq x) - \Phi(x)| \leq \frac{C\rho}{\sigma^3\sqrt n}$ with an absolute constant $C < 0.48$; the rate $n^{-1/2}$ is attained by
Bernoulli variables \cite{durrett-probability}.
\item (\emph{Feller}) If $\max_m\sigma_{n,m}^2 \to 0$, Lindeberg's condition is also necessary for $S_n \Rightarrow N(0, \sigma^2)$.
\item (\emph{Law of the iterated logarithm}) For i.i.d.\ variables with mean $0$ and variance $1$, $\limsup_n\frac{S_n}{\sqrt{2n\log\log n}} = 1$ almost surely (Hartman--Wintner): the central limit theorem says $S_n/\sqrt n$ is of order $1$
in distribution, but along a typical path it fluctuates up to $\sqrt{2\log\log n}$ \cite{durrett-probability}.
\item The multivariate central limit theorem holds with the same proof, using $\varphi_X(t) = E[e^{i\langle t, X\rangle}]$ on $\RR^d$: i.i.d.\ random vectors with covariance matrix $\Sigma$ give a Gaussian limit with covariance $\Sigma$.
\end{enumerate}
\end{remark}

\begin{counterexample}\label{probability:counterexample-clt}
Finite variance is essential: for i.i.d.\ Cauchy variables, $S_n/\sqrt n = \sqrt n\cdot S_n/n$ has the law of $\sqrt n$ times a Cauchy variable, which does not converge in distribution; the correct normalisation is $S_n/n$, and the
limit is Cauchy, not normal. Heavy-tailed variables lead to the stable laws of Section~\ref{probability:section-poisson}.
\end{counterexample}

\section{Poisson approximation, infinitely divisible and stable laws}\label{probability:section-poisson}

\begin{theorem}[Law of rare events]\label{probability:theorem-poisson-limit}
If $S_n \sim B(n, p_n)$ with $np_n \to \lambda > 0$, then $S_n \Rightarrow \mathrm{Poi}(\lambda)$; equivalently, $P(S_n = k) \to e^{-\lambda}\lambda^k/k!$ for every $k$.
\end{theorem}

\begin{proof}
$\varphi_{S_n}(t) = (1 + p_n(e^{it} - 1))^n = (1 + \frac{z_n}{n})^n$ with $z_n = np_n(e^{it} - 1) \to z = \lambda(e^{it} - 1)$, and $(1 + z_n/n)^n \to e^z$ (take logarithms for large $n$). L\'evy's theorem and
Example~\ref{probability:example-characteristic-functions}(1) give $S_n \Rightarrow \mathrm{Poi}(\lambda)$; for integer-valued variables, narrow convergence gives convergence of each $P(S_n = k)$, by testing against a continuous
function equal to $1$ at $k$ and $0$ outside $(k - \frac12, k + \frac12)$.
\end{proof}

\begin{definition}\label{probability:definition-total-variation}
The \emph{total variation distance} between laws $\mu$, $\nu$ on a countable set is $d_{TV}(\mu, \nu) = \frac12\sum_k|\mu(k) - \nu(k)| = \sup_B|\mu(B) - \nu(B)|$.
\end{definition}

\begin{proposition}[Le Cam]\label{probability:proposition-le-cam}
Let $X_1, \ldots, X_n$ be independent Bernoulli variables with $P(X_i = 1) = p_i$, and $S = \sum X_i$. Then $d_{TV}(P_S, \mathrm{Poi}(\sum p_i)) \leq \sum p_i^2$.
\end{proposition}

\begin{proof}
\emph{Coupling.} For laws $\mu$, $\nu$ on $\ZZ_{\geq 0}$ there are random variables $X \sim \mu$, $Y \sim \nu$ on one probability space with $P(X \neq Y) = d_{TV}(\mu, \nu)$: let $m(k) = \min(\mu(k), \nu(k))$, $\alpha = \sum m = 1 - d_{TV}$; with
probability $\alpha$ choose $X = Y$ according to $m/\alpha$, and otherwise choose $X$ according to $(\mu - m)/(1 - \alpha)$ and $Y$ according to $(\nu - m)/(1 - \alpha)$ (these have disjoint supports). \emph{One summand.} For
$\mu = B(1, p)$ and $\nu = \mathrm{Poi}(p)$, $2d_{TV} = (e^{-p} - 1 + p) + p(1 - e^{-p}) + (1 - e^{-p} - pe^{-p}) = 2p(1 - e^{-p}) \leq 2p^2$. \emph{Sum.} Couple each $X_i$ with $Y_i \sim \mathrm{Poi}(p_i)$, independently over $i$
(Corollary~\ref{measure-theory:corollary-independent-sequences}). Then $\sum Y_i \sim \mathrm{Poi}(\sum p_i)$ (Example~\ref{probability:example-sums}) and
$d_{TV}(P_S, P_{\sum Y_i}) \leq P(\sum X_i \neq \sum Y_i) \leq \sum P(X_i \neq Y_i) \leq \sum p_i^2$.
\end{proof}

So the number of rare, roughly independent events is approximately Poisson: misprints on a page, decays of atoms in a second, Prussian cavalrymen killed by horse kicks (Bortkiewicz's famous data of 1898).

\begin{definition}\label{probability:definition-infinitely-divisible}
A law $\mu$ on $\RR$ is \emph{infinitely divisible} if for every $n$ there is a law $\mu_n$ with $\mu = \mu_n^{*n}$, that is, $X \overset{d}{=} X_{n,1} + \cdots + X_{n,n}$ with i.i.d.\ $X_{n,m}$; equivalently, $\varphi_\mu$ has an $n$-th root
that is a characteristic function, for every $n$. It is \emph{stable} if for i.i.d.\ $X$, $X_1, X_2, \ldots \sim \mu$ and every $n$ there are $a_n > 0$, $b_n$ with $X_1 + \cdots + X_n \overset{d}{=} a_nX + b_n$.
\end{definition}

\begin{example}\label{probability:example-infinitely-divisible}
Normal, Poisson, Cauchy and Gamma laws are infinitely divisible: $e^{imt - \sigma^2t^2/2}$, $e^{\lambda(e^{it} - 1)}$, $e^{-|t|}$ and $(1 - it/\lambda)^{-\alpha}$ have the $n$-th roots $e^{imt/n - \sigma^2t^2/2n}$, etc., of the same type. So is a
\emph{compound Poisson} law, the law of $Y_1 + \cdots + Y_N$ with $N \sim \mathrm{Poi}(\lambda)$ independent of i.i.d.\ $Y_k \sim \nu$, whose characteristic function is $\exp(\lambda\int(e^{itx} - 1)\,d\nu(x))$. Normal and Cauchy laws are
stable, with $a_n = \sqrt n$ and $a_n = n$.
\end{example}

\begin{counterexample}\label{probability:counterexample-not-infinitely-divisible}
A nondegenerate law with bounded support is never infinitely divisible: if $|X| \leq c$ and $X \overset{d}{=} \sum_{m \leq n}X_{n,m}$, then $|X_{n,m}| \leq c/n$ almost surely (if $P(X_{n,1} > c/n) > 0$, all $n$ summands exceed $c/n$ simultaneously with positive probability, and the sum exceeds $c$; similarly below), so
$\operatorname{Var}X = n\operatorname{Var}X_{n,1} \leq c^2/n \to 0$. In particular Bernoulli and uniform laws are not infinitely divisible. Also $e^{-|t|^\alpha}$ is not a characteristic function for $\alpha > 2$: it would be twice
differentiable at $0$ with second derivative $0$, but $E[X^2] \leq \liminf_{h \to 0}\frac{2 - \varphi(h) - \varphi(-h)}{h^2} = -\varphi^{\prime\prime}(0) = 0$ by Fatou's lemma applied to $\frac{2 - 2\cos hX}{h^2} \to X^2$, so $X = 0$ and
$\varphi = 1$, a contradiction.
\end{counterexample}

\begin{theorem}[L\'evy--Khintchine]\label{probability:theorem-levy-khintchine}
A function $\varphi$ is the characteristic function of an infinitely divisible law iff
\[
\varphi(t) = \exp\Bigl(ibt - \frac{\sigma^2t^2}{2} + \int_{\RR \setminus 0}\bigl(e^{itx} - 1 - itx\mathbf{1}_{\{|x| \leq 1\}}\bigr)\,d\nu(x)\Bigr)
\]
with $b \in \RR$, $\sigma^2 \geq 0$ and a measure $\nu$ on $\RR \setminus 0$ (the \emph{L\'evy measure}) with $\int\min(1, x^2)\,d\nu < \infty$. The stable laws are the normal laws and, for $0 < \alpha < 2$, the laws with
$\nu(dx) = (c_+\mathbf{1}_{x > 0} + c_-\mathbf{1}_{x < 0})|x|^{-1-\alpha}dx$; the symmetric ones have $\varphi(t) = e^{-c|t|^\alpha}$, $0 < \alpha \leq 2$.
\end{theorem}

This is quoted from \cite[Chapter XVII]{feller-2}. Infinitely divisible laws are exactly the limits of row sums of triangular arrays of independent, uniformly small variables, and the laws of $L_1$ for L\'evy processes
$(L_t)$ (Chapter~\ref{martingales}); $\sigma^2$ describes the Brownian part and $\nu$ the jumps.

\section{Large deviations}\label{probability:section-large-deviations}

The law of large numbers says that $P(S_n \geq nx) \to 0$ for $x > m$; the central limit theorem describes deviations of size $\sqrt n$. Deviations of size $n$ have exponentially small probability, with an exact rate.

\begin{definition}\label{probability:definition-rate-function}
For a real random variable $X$ with $\Lambda(\lambda) = \log E[e^{\lambda X}] < \infty$ for all $\lambda$, the \emph{rate function} is the Legendre transform $I(x) = \sup_{\lambda \in \RR}(\lambda x - \Lambda(\lambda))$.
\end{definition}

\begin{theorem}[Cram\'er]\label{probability:theorem-cramer}
Let $X_1, X_2, \ldots$ be i.i.d.\ copies of $X$, with $\Lambda$ finite everywhere and mean $m$, and let $x > m$ be such that $\Lambda^{\prime}(\lambda^*) = x$ for some $\lambda^* > 0$. Then
\[
\lim_{n \to \infty}\frac1n\log P(S_n \geq nx) = -I(x) = -(\lambda^*x - \Lambda(\lambda^*)).
\]
\end{theorem}

\begin{proof}
$\Lambda$ is smooth and convex (differentiate under the integral; $\Lambda^{\prime\prime}$ is the variance of the tilted law below), so $\lambda \mapsto \lambda x - \Lambda(\lambda)$ is concave with critical point $\lambda^*$, and
$I(x) = \lambda^*x - \Lambda(\lambda^*)$. \emph{Upper bound.} By the Chernoff bound (Proposition~\ref{probability:proposition-chernoff}) with $\lambda = \lambda^*$, $P(S_n \geq nx) \leq e^{-nI(x)}$ for every $n$.
\emph{Lower bound by tilting.} Let $Q$ be the law with $dQ = e^{\lambda^*y - \Lambda(\lambda^*)}dP_X(y)$, a probability measure with mean $\Lambda^{\prime}(\lambda^*) = x$ and finite variance $\Lambda^{\prime\prime}(\lambda^*)$. Let $Y_k$ be
i.i.d.\ $\sim Q$ and $T_n = \sum Y_k$; then for every Borel $B \subseteq \RR^n$, $P((X_1, \ldots, X_n) \in B) = E[e^{-\lambda^*T_n + n\Lambda(\lambda^*)}; (Y_1, \ldots, Y_n) \in B]$. For $\delta > 0$,
\[
P(S_n \geq nx) \geq P(nx \leq S_n \leq n(x + \delta)) \geq e^{-n(\lambda^*(x + \delta) - \Lambda(\lambda^*))}P(0 \leq T_n - nx \leq n\delta).
\]
By the central limit theorem for the $Y_k$, $P(0 \leq T_n - nx \leq n\delta) \to \frac12$. Hence $\liminf\frac1n\log P(S_n \geq nx) \geq -I(x) - \lambda^*\delta$ for every $\delta > 0$.
\end{proof}

\begin{example}\label{probability:example-cramer}
\begin{enumerate}
\item For $X \sim N(0, 1)$, $\Lambda(\lambda) = \lambda^2/2$ and $I(x) = x^2/2$: $P(S_n \geq nx) = e^{-nx^2/2 + o(n)}$, consistent with $S_n/\sqrt n \sim N(0, 1)$.
\item For a fair coin, $X \in \{0, 1\}$, $\Lambda(\lambda) = \log\frac{1 + e^\lambda}{2}$ and $I(x) = x\log 2x + (1 - x)\log 2(1 - x)$ for $0 < x < 1$, the \emph{relative entropy} of $(x, 1 - x)$ with respect to $(\frac12, \frac12)$. The
probability of at least $60\%$ heads in $n$ tosses is $e^{-0.0201n + o(n)}$.
\end{enumerate}
\end{example}

\begin{remark}\label{probability:remark-sanov}
Cram\'er's theorem holds for all $x$ (with $I = \infty$ outside the closed convex hull of the support), and without assuming $\Lambda$ finite, in the general form of a \emph{large deviation principle}. Sanov's theorem is
the analogue for empirical measures: the probability that the empirical distribution of $n$ i.i.d.\ samples from $\mu$ is close to $\nu$ is $e^{-nH(\nu \mid \mu) + o(n)}$, with $H(\nu \mid \mu) = \int\log\frac{d\nu}{d\mu}\,d\nu$ the
relative entropy \cite{dembo-zeitouni}. This connects probability with information theory (Chapter~\ref{information-theory}) and with statistical mechanics, where $-I$ is an entropy.
\end{remark}

\section{Random matrices}\label{probability:section-random-matrices}

\begin{definition}\label{probability:definition-wigner}
A \emph{Wigner matrix} of size $n$ is a real symmetric random matrix $W_n = (w_{ij})$ whose entries $w_{ij}$, $i \leq j$, are independent, with $w_{ij}$ ($i < j$) identically distributed with mean $0$ and variance $1$, the diagonal
entries identically distributed with mean $0$, and all moments finite. Its \emph{empirical spectral distribution} is $L_n = \frac1n\sum_{k=1}^n\delta_{\lambda_k/\sqrt n}$, where $\lambda_1, \ldots, \lambda_n$ are the eigenvalues of $W_n$.
The \emph{semicircle law} is $\sigma(dx) = \frac{1}{2\pi}\sqrt{4 - x^2}\,\mathbf{1}_{[-2,2]}(x)\,dx$.
\end{definition}

\begin{proposition}\label{probability:proposition-semicircle-moments}
The semicircle law is a probability measure; its odd moments vanish, and $\int x^{2k}\,d\sigma = C_k = \frac{1}{k + 1}\binom{2k}{k}$, the $k$-th Catalan number.
\end{proposition}

\begin{proof}
Odd moments vanish by symmetry. Substituting $x = 2\sqrt s$ on $[0, 2]$,
$\int x^{2k}d\sigma = \frac1\pi\int_0^2x^{2k}\sqrt{4 - x^2}\,dx = \frac{2\cdot 4^k}{\pi}\int_0^1s^{k-1/2}(1 - s)^{1/2}ds = \frac{2\cdot 4^k}{\pi}B(k + \tfrac12, \tfrac32)$. By Proposition~\ref{special-functions:proposition-beta}, with
$\Gamma(\frac32) = \frac{\sqrt\pi}{2}$ and $\Gamma(k + \frac12) = \frac{(2k)!}{4^kk!}\sqrt\pi$ (from $\Gamma(\frac12) = \sqrt\pi$ and the functional equation), this is
$\frac{2\cdot 4^k}{\pi}\cdot\frac{(2k)!\sqrt\pi}{4^kk!}\cdot\frac{\sqrt\pi}{2}\cdot\frac{1}{(k + 1)!} = \frac{(2k)!}{k!(k + 1)!}$. For $k = 0$ this gives total mass $1$.
\end{proof}

\begin{theorem}[Wigner's semicircle law]\label{probability:theorem-wigner}
For Wigner matrices, $E\int x^k\,dL_n \to \int x^k\,d\sigma$ for every $k$, and $L_n \to \sigma$ narrowly in probability: for every bounded continuous $h$, $\int h\,dL_n \to \int h\,d\sigma$ in probability.
\end{theorem}

\begin{proof}[Sketch of proof]
\emph{Moments as walks.} $\int x^k\,dL_n = \frac{1}{n^{1+k/2}}\operatorname{tr}W_n^k = \frac{1}{n^{1+k/2}}\sum_{i_1, \ldots, i_k}w_{i_1i_2}w_{i_2i_3}\cdots w_{i_ki_1}$, a sum over closed walks of length $k$ on $\{1, \ldots, n\}$. \emph{Counting.} The
expectation of a term vanishes unless every edge is traversed at least twice (mean $0$ and independence). A walk with $k$ steps traversing each of its edges at least twice visits at most $\frac k2 + 1$ vertices, with equality iff
$k$ is even, each edge is traversed exactly twice and the edges form a tree; these walks are counted by $n^{k/2+1}(1 + o(1))$ times the number of rooted plane trees with $k/2$ edges, which is $C_{k/2}$ (by
Theorem~\ref{generating-functions:theorem-catalan} and the interpretations of Proposition~\ref{generating-functions:proposition-catalan-interpretations}), each contributing $E[w_{12}^2]^{k/2} = 1$. All other walks contribute $O(n^{-1/2})$
after normalisation, using the finite moments. \emph{Concentration.} A similar count shows $\operatorname{Var}\int x^k\,dL_n = O(n^{-2})$. \emph{Moments to measures.} $\sigma$ has compact support, so it is determined by its moments,
and moment convergence plus concentration gives narrow convergence in probability. The omitted verifications are the error terms in the counting \cite{anderson-guionnet-zeitouni}.
\end{proof}

\begin{remark}\label{probability:remark-random-matrices}
For Gaussian entries (the \emph{Gaussian orthogonal ensemble}, or its complex Hermitian version, the GUE) the joint density of the eigenvalues is explicit, proportional to $\prod_{i<j}|\lambda_i - \lambda_j|^\beta e^{-c\sum\lambda_i^2}$
with $\beta = 1$ or $2$ and a constant $c > 0$ depending on the normalisation: eigenvalues repel each other. The largest eigenvalue, rescaled as $n^{2/3}(\lambda_{\max}/\sqrt n - 2)$, converges to the \emph{Tracy--Widom distribution}, which also describes the longest increasing
subsequence of a random permutation (Baik--Deift--Johansson). For sample covariance matrices $\frac1nXX^T$ with $X$ of size $p \times n$, $p/n \to c$, the limit is the \emph{Marchenko--Pastur law}. Montgomery and Dyson observed that
the zeros of the Riemann zeta function, suitably rescaled, have the pair correlation of GUE eigenvalues (Chapter~\ref{analytic-number-theory}) \cite{anderson-guionnet-zeitouni}.
\end{remark}

\section{Exercises}\label{probability:section-exercises}

\begin{exercise}\label{probability:exercise-coupon-collector}
A collector buys cards, each uniformly random among $n$ types, independently. Let $T_n$ be the number of purchases until all types are collected. Show that $E[T_n] = nH_n$, $H_n = 1 + \frac12 + \cdots + \frac1n$, that
$\operatorname{Var}T_n \leq \frac{\pi^2}{6}n^2$, and that $T_n/(n\log n) \to 1$ in probability.
\end{exercise}

\begin{solution}
After $k - 1$ types are collected, the waiting time for a new type is geometric with success probability $p_k = \frac{n - k + 1}{n}$, and these waiting times $G_k$ are independent; $T_n = \sum_{k=1}^nG_k$. So
$E[T_n] = \sum\frac{n}{n - k + 1} = nH_n$ and $\operatorname{Var}T_n = \sum\frac{1 - p_k}{p_k^2} \leq \sum_{j=1}^n\frac{n^2}{j^2} \leq \frac{\pi^2}{6}n^2$. By Chebyshev, $P(|T_n - nH_n| \geq \varepsilon n\log n) \leq \frac{\pi^2}{6\varepsilon^2\log^2n} \to 0$, and
$H_n/\log n \to 1$.
\end{solution}

\begin{exercise}\label{probability:exercise-memoryless}
Show that a random variable $X > 0$ with $P(X > s + t) = P(X > s)P(X > t)$ for all $s, t \geq 0$ is exponential (or $X = \infty$ a.s.).
\end{exercise}

\begin{solution}
$g(t) = P(X > t)$ is nonincreasing, right-continuous, with $g(s + t) = g(s)g(t)$. If $g(1) = 0$ then $g(1/n)^n = 0$ and $g(1/n) = 0$, so $P(X > 0) = \lim g(1/n) = 0$, impossible. So $g(1) = e^{-\lambda}$ with $\lambda \in [0, \infty)$, then
$g(q) = e^{-\lambda q}$ for rational $q > 0$ by the functional equation, and for all $t$ by monotonicity. If $\lambda = 0$, $X = \infty$ almost surely; otherwise $X \sim \mathrm{Exp}(\lambda)$.
\end{solution}

\begin{exercise}\label{probability:exercise-birthday}
Show that among $k$ people with independent uniformly distributed birthdays among $N$ days, the probability that all birthdays differ is at most $e^{-k(k-1)/2N}$; deduce that for $N = 365$ and $k = 23$ a coincidence is more likely
than not.
\end{exercise}

\begin{solution}
The probability is $\prod_{j=1}^{k-1}(1 - \frac jN) \leq \prod e^{-j/N} = e^{-k(k-1)/2N}$, using $1 - x \leq e^{-x}$. For $k = 23$, $N = 365$: $e^{-253/365} = e^{-0.693\ldots} < \frac12$, since $\log 2 = 0.6931\ldots$ and $253/365 = 0.6932\ldots$.
\end{solution}

\begin{exercise}\label{probability:exercise-characteristic-symmetrisation}
Show that if $\varphi$ is a characteristic function, so are $\overline\varphi$, $|\varphi|^2$ and $\operatorname{Re}\varphi$. Show that $\cos t$ is the characteristic function of a fair sign $\pm1$, but $\cos t$ is not the characteristic function
of an infinitely divisible law.
\end{exercise}

\begin{solution}
$\overline\varphi = \varphi_{-X}$, $|\varphi|^2 = \varphi_{X - X^{\prime}}$ for an independent copy $X^{\prime}$, and $\operatorname{Re}\varphi$ is the characteristic function of $\varepsilon X$ with an independent fair sign $\varepsilon$, as it equals
$\frac12(\varphi + \overline\varphi)$. $E[e^{it\varepsilon}] = \frac12(e^{it} + e^{-it}) = \cos t$. The sign has bounded support and is nondegenerate, so it is not infinitely divisible
(Counterexample~\ref{probability:counterexample-not-infinitely-divisible}).
\end{solution}

\begin{exercise}\label{probability:exercise-chernoff-coin}
Using Proposition~\ref{probability:proposition-chernoff}, show that for $n$ fair coin tosses $P(S_n \geq \frac34n) \leq (2\cdot 3^{-3/4})^n \approx 0.877^n$.
\end{exercise}

\begin{solution}
$M(\lambda) = \frac{1 + e^\lambda}{2}$. Minimise $e^{-3\lambda/4}\frac{1 + e^\lambda}{2}$: the derivative vanishes at $e^\lambda = 3$, giving $3^{-3/4}\cdot\frac{1 + 3}{2} = 2\cdot 3^{-3/4} = 0.877\ldots$. (The exponent is $I(\frac34) = \frac34\log\frac32 + \frac14\log\frac12 = 0.1308\ldots$, and
$e^{-0.1308} = 0.877$.)
\end{solution}

\begin{exercise}\label{probability:exercise-poisson-thinning}
Let $N \sim \mathrm{Poi}(\lambda)$ and, given $N$, let each of the $N$ points be kept independently with probability $p$. Show that the number of kept points is $\mathrm{Poi}(\lambda p)$, independent of the number of discarded
points, which is $\mathrm{Poi}(\lambda(1 - p))$.
\end{exercise}

\begin{solution}
$P(K = k, D = d) = P(N = k + d)\binom{k + d}{k}p^k(1 - p)^d = e^{-\lambda}\frac{\lambda^{k+d}}{k!\,d!}p^k(1 - p)^d = e^{-\lambda p}\frac{(\lambda p)^k}{k!}\cdot e^{-\lambda(1-p)}\frac{(\lambda(1 - p))^d}{d!}$, a product of the two Poisson mass functions.
\end{solution}

\begin{exercise}\label{probability:exercise-clt-gamma}
Deduce from the central limit theorem that $e^{-n}\sum_{k=0}^n\frac{n^k}{k!} \to \frac12$.
\end{exercise}

\begin{solution}
The sum is $P(S_n \leq n)$ for $S_n \sim \mathrm{Poi}(n)$, a sum of $n$ independent $\mathrm{Poi}(1)$ variables with mean and variance $1$ (Example~\ref{probability:example-sums}). By Corollary~\ref{probability:corollary-clt},
$P(\frac{S_n - n}{\sqrt n} \leq 0) \to \Phi(0) = \frac12$.
\end{solution}
